\documentclass[11pt,a4paper]{article}
\usepackage[T1]{fontenc}
\usepackage[utf8]{inputenc}
\usepackage[french]{babel}
\usepackage{lmodern}
\usepackage{amsmath,amssymb}
\usepackage[margin=2.6cm]{geometry}
\usepackage{booktabs}
\usepackage{enumitem}
\usepackage{xcolor}
\usepackage{hyperref}
\hypersetup{colorlinks=true,linkcolor=blue!50!black}

\newcommand{\NW}{\mathrm{NW}}
\newcommand{\code}[1]{\texttt{#1}}

\title{Rapport d'implémentation M3\\
\large Livrable 03 du programme M3 (workspace \code{m3\_credit\_soc})}
\author{Agent de recherche M3}
\date{4 juillet 2026}

\begin{document}
\maketitle

\begin{abstract}
Le modèle M3 spécifié dans le livrable 02 est implémenté dans
\code{src/m3/} (11 modules, $\sim$1500 lignes), couvert par 50 tests
(mécaniques, contrats et marché, faillites et avalanches causales,
invariants I1--I9, propriété de réduction R1, estimateurs statistiques sur
données synthétiques), tous verts. La propriété de réduction R1 --- M3 avec
$s{=}1$, $c{=}0$, $L_0{=}0$ et crédit coupé reproduit \code{m2\_fable} bit à
bit à seed égale --- est vérifiée exactement, ce qui ancre M3 sur le modèle
nul M2 sans dérive d'implémentation. La calibration pré-enregistrée a figé
$(s, c) = (0{,}75\,;\,0{,}10)$ ; la cellule provisoire $(0{,}75\,;\,0{,}05)$
échoue au critère de régime --- résultat mécanistique en soi (le coussin de
liquidité éteint la mortalité). Un bug d'observation non neutre (violation
de I8 par mutation de \code{defaultdict}) a été détecté par les tests et
corrigé.
\end{abstract}

\section{Architecture}

\begin{center}
\small
\begin{tabular}{lp{9.6cm}}
\toprule
Module & Rôle \\
\midrule
\code{config.py} & \code{M3Config} gelée (dataclass \code{frozen}) ; toute
variante expérimentale est un champ nommé ; assertions de domaine. \\
\code{entities.py} & Population en structure de tableaux ($L$, $K$, vivante,
naissance, secteur, flux du pas, rente E) ; ids jamais réutilisés. \\
\code{contracts.py} & \code{LoanBook} : contrats nominaux perpétuels ;
agrégats \code{claims}/\code{debts}/\code{due} maintenus incrémentalement
(I2/I3) ; fusion par paire (\code{merge\_pairs}) au taux moyen pondéré ;
\code{check\_consistency} de contrôle hors boucle. \\
\code{production.py} & Choc multiplicatif (i.i.d. ou corrélé macro/sectoriel,
chaque composante centrée à $-\mathrm{var}/2$), production--partage
($K \mathrel{+}= sP$, $L \mathrel{+}= (1{-}s)P$), dépréciation et
consommation. \\
\code{market.py} & Marché local : liquidité prêtable
$A = \max(0, L - \beta_L \cdot \mathrm{due})$, sélection assortative
($\arg\max A$ / $\arg\min K$) ou aléatoire (ablation F), taux géométrique,
cible $K^*(\rho)$, transfert $L_\ell \to K_b$ (ou $L_b$, ablation C). \\
\code{bankruptcy.py} & Cascade au point fixe (I9) ; pertes de stock et de
flux ; variantes D (indemnisation) et E (rente fantôme) ; arêtes causales de
perte ; avalanches = composantes connexes du graphe de pertes entre faillies
du même pas (union-find), profondeur = itération maximale, causes des
racines (liquidité / insolvabilité / les deux). \\
\code{simulation.py} & Boucle à 8 phases dans l'ordre pré-enregistré ;
séries par pas (27 champs, dont bilan I4 : gain de choc, déprécié, consommé,
détruit, injecté) ; instantanés individus + réseau ; RNG local
\code{default\_rng(seed)} unique. \\
\code{metrics.py} & Écriture/lecture des artefacts : \code{series.csv},
\code{snap\_t*.npz}, \code{net\_t*.npz}, \code{avalanches.csv},
\code{config.json}, \code{summary.json} (avec \code{wall\_seconds} pour le
budget). \\
\code{analysis.py} & Outillage statistique corrigé hérité de
\code{m2\_fable} (MLE tronqués, LR renormalisé, $x_{\min}$ commun,
bootstrap) + familles complètes (dont dPlN), Gini, HHI, parts du top,
matrices de transition, AUC réseau$\to$défauts, distribution des tailles
d'avalanches causales. \\
\code{dpln.py} & Double Pareto-lognormale : log-pdf Normal-Laplace stable
(\code{logcdf} + \code{logaddexp}), générateur, MLE multi-départs. \\
\code{plots.py} & Figures générées depuis les fichiers de résultats
uniquement. \\
\bottomrule
\end{tabular}
\end{center}

Expériences sous \code{experiments/m3/} : \code{run\_calibration.py}
(pré-enregistrée), \code{run\_baseline.py}, \code{run\_ablation.py} (B--J),
\code{run\_grid.py}, \code{run\_long.py}, \code{run\_validation.py}
(analyse d'un run $\to$ \code{validation.json}, jamais de pooling
inter-seed), \code{aggregate\_results.py} (table comparative des signatures),
\code{exp\_common.py} (exécution + comptabilité du budget 48 h).

\section{Tests (50, tous verts)}

\begin{itemize}[leftmargin=1.5em, itemsep=2pt]
\item \textbf{Mécaniques} (8) : naissance et $\NW$ implicite, partage exact
  de la production, absence de drift du choc (i.i.d., macro, sectoriel —
  composante commune vérifiée à variance idiosyncratique nulle),
  dépréciation/consommation exactes, non-dépréciation du nominal, mort.
\item \textbf{Contrats et marché} (8) : agrégats incrémentaux, fusion par
  paire (flux d'intérêts exactement préservé), I5, formules $q$, $r$,
  $K^*(\rho)$ vérifiées contre l'implémentation, ablation C, tampon
  $\beta_L$, marché aléatoire, \code{credit=False} inerte.
\item \textbf{Faillites et cascades} (14) : service complet/partiel prorata,
  défaut de liquidité d'une entité \emph{solvable} (cause « liquidity »),
  pertes de stock et de flux, résidu en nature prorata/détruit, cascade à
  deux niveaux (une avalanche, taille 2, profondeur 2, une racine), racines
  indépendantes non agrégées, ablations D et E (rente versée au pas suivant,
  éteinte à la mort), transfert/annulation des prêts de la faillie, aucun
  contrat orphelin.
\item \textbf{Invariants} (6) : I1 ; I4 --- bilan global vérifié pas à pas à
  $10^{-6}$ relatif sur un run avec crédit ; I7 ; I8 strict (bit à bit) ;
  I9 ; R1.
\item \textbf{Estimateurs} (14) : chaque estimateur validé sur données
  synthétiques à vérité connue avant tout usage réel --- dont le LR corrigé
  ($R<0$ sur lognormale pure ; $|z|<2$ sur Pareto pure : test peu puissant,
  documenté) et le MLE dPlN (recouvrement de $a$ et $b$).
\end{itemize}

\paragraph{La propriété de réduction R1, ancre anti-dérive.} M3 avec
$s{=}1$, $c{=}0$, $L_0{=}0$, crédit coupé émet exactement la même séquence
RNG et les mêmes opérations flottantes que \code{m2\_fable}
(\code{credit=False}) : à $T{=}200$ et $T{=}300$, $\max_i |K_i - w_i| = 0.0$,
vivantes et morts par pas identiques. Toute conclusion « B diffère de M2 »
sera donc attribuable aux mécanismes nouveaux ($s<1$, $c>0$, $L_0>0$), pas à
une divergence d'implémentation.

\section{Bug détecté par les tests : observation non neutre (I8)}

La version initiale de \code{snapshot()} lisait \code{book.claims[i]},
\code{book.by\_borrower[i]}, etc. Or ces structures sont des
\code{defaultdict} : l'accès par \code{[]} \emph{insère} une clé vide pour
chaque vivante interrogée. La phase de service des intérêts itérant sur
\code{book.by\_borrower.keys()} (ordre d'insertion), un run \emph{avec}
instantanés servait les emprunteuses dans un autre ordre qu'un run
\emph{sans} --- sommes flottantes réordonnées, divergence au dernier ulp dès
$t{=}54$ (seed 9), potentiellement amplifiable par un seuil de faillite sur
un run long. Correctif : accès en \code{.get()} (lecture pure) ; le test I8
strict (identité bit à bit des séries) est rétabli et vert. Leçon
générale consignée dans \code{NOTES.md} : toute fonction d'observation doit
être en lecture pure ; sur un \code{defaultdict}, \code{[]} est une
écriture.

\section{Calibration pré-enregistrée : $(s, c)$ figés à $(0{,}75\,;\,0{,}10)$}

Scan B (sans crédit) de $s \in \{0{,}5, 0{,}65, 0{,}75, 0{,}9\} \times
c \in \{0{,}02, 0{,}05, 0{,}10\}$, $T{=}1500$, seed 0, critère du rapport 02
§4. Résultat : \textbf{seule la colonne $c = 0{,}10$ est viable} (les quatre
cellules ; croissance du dernier quart $\le 12\,\%$, morts/naissances
$0{,}94$--$1{,}01$). À $c \le 0{,}05$, la population diverge
(morts/naissances $0{,}34$--$0{,}44$, $N > 8000$ et croissant) : le stock
stationnaire de liquidité $L^* \approx (1{-}s)\sqrt{K^*}/c$ forme un coussin
au-dessus du plancher $d_0$ qui éteint la mortalité endogène. La cellule
provisoire $(0{,}75\,;\,0{,}05)$ échoue donc ; la règle pré-enregistrée
désigne $(0{,}75\,;\,0{,}10)$, figée dans \code{config.py} et non retouchée
ensuite. Enseignement mécanistique (pour le rapport 06) : la mortalité de M2
repose sur l'absence de tout actif stable amortisseur ; tout mécanisme
ajoutant un stock sûr à $\NW$ referme la fenêtre démographique --- l'axe $c$
de la grille est le levier de régime dominant, plus que $s$.

\section{Limites numériques et performance}

\begin{itemize}[leftmargin=1.5em, itemsep=2pt]
\item Gardes : $q_{\min} = 10^{-9}$ (contrats), \code{w\_clamp} $= 10^{-12}$
  (arrondi à zéro de $L$, $K$), \code{tol\_nw} $= 10^{-9}$ (faillite),
  \code{pop\_max} $= 30\,000$ (arrêt, ne borne pas la dynamique),
  plancher $10^{-9}$ dans les taux marginaux.
\item \code{merge\_pairs} actif par défaut : carnet borné par le nombre de
  paires actives (leçon M2 : le prorata pur explose combinatoirement).
\item La rente fantôme (E) et l'indemnisation (D) sont des injections
  \emph{délibérées}, mesurées dans la série (champ \code{injected}) et dans
  le bilan I4.
\item Coûts mesurés (1 cœur) : run B $T{=}1500$ : 10--30 s ; run A (crédit
  actif) $T{=}600$, $N \approx 1100$ : 8 s. Boucles en listes Python comme
  M2 (choix assumé : lisibilité et identité R1 plutôt que vectorisation).
\item Suivi du budget 48 h : somme des \code{wall\_seconds} des
  \code{summary.json} (\code{exp\_common.py}), journalisée dans
  \code{NOTES.md}.
\end{itemize}

\section{Différences avec M2 (exhaustives)}

(1) état $(L, K)$ au lieu de $w$ scalaire ; (2) production partagée
$s$/$1{-}s$ au lieu d'entièrement retenue ; (3) consommation $cL$ (nouveau
puits) ; (4) intérêts exigibles en liquidité seulement, service simultané au
prorata par emprunteuse (M2 \code{fable} : ordre de création — convention
codex adoptée, documentée) ; (5) prêt $L_\ell \to K_b$ au lieu de
$w_\ell \to w_b$ ; (6) liquidité prêtable bornée par le tampon $\beta_L$ ;
(7) résidu de faillite réparti en nature ($L$ et $K$ séparément) ;
(8) défaut de liquidité distinct de l'insolvabilité, tracé par cause ;
(9) avalanches causales (graphe de pertes) au lieu de lots par pas ;
(10) chocs corrélables (G) ; (11) variantes D/E/F nouvelles ;
(12) \code{merge\_pairs} par défaut (opt-in dans M2).
Tout le reste (naissances Poisson, $d_0$, choc sans drift, extraction
concave, dépréciation réelle, taux géométrique, cible $K^*(\rho)$,
$k$-échantillonnage, cascade au point fixe, conventions de résidu) est
conservé à l'identique --- garanti par R1.

\end{document}
